\documentclass[11pt]{article}
\usepackage[T1]{fontenc}
\usepackage[utf8]{inputenc}
\usepackage{lmodern,amsmath,amssymb,amsthm,mathtools,mathrsfs}
\usepackage[margin=25mm]{geometry}
\usepackage{booktabs,tabularx,longtable,array,enumitem}
\usepackage{xcolor,hyperref,microtype}
\hypersetup{colorlinks=true,linkcolor=blue!45!black,citecolor=blue!45!black,urlcolor=blue!45!black}
\setlength{\emergencystretch}{3em}
\setlist{nosep,leftmargin=*}
\newcommand{\Ric}{\operatorname{Ric}}
\newcommand{\Rm}{\operatorname{Rm}}
\newcommand{\Vol}{\operatorname{Vol}}
\newcommand{\tr}{\operatorname{tr}}
\newcommand{\diver}{\operatorname{div}}
\newcommand{\Id}{\operatorname{Id}}
\newcommand{\Met}{\operatorname{Met}}
\newcommand{\CP}{\mathbb{CP}}
\newcommand{\R}{\mathbb R}
\newcommand{\C}{\mathbb C}
\newcommand{\dd}{\,dV_g}
\newcommand{\cE}{\mathcal E}
\newcommand{\cJ}{\mathcal J}
\newcommand{\pstar}{\mathfrak p_*}
\newcommand{\src}[1]{\textsf{[#1]}}
\newtheorem{theorem}{Theorem}[section]
\newtheorem{proposition}[theorem]{Proposition}
\newtheorem{lemma}[theorem]{Lemma}
\newtheorem{corollary}[theorem]{Corollary}
\newtheorem{question}[theorem]{Question}
\theoremstyle{definition}
\newtheorem{definition}[theorem]{Definition}
\theoremstyle{remark}
\newtheorem{remark}[theorem]{Remark}
\title{Einstein Gaps on Four-Manifolds\\[4pt]
\large Subsequent Mathematical Developments:\\
Explicit Almost-K\"ahler Geometry, Harmonic Flux,\\
Variational Descent, and Global Upper Bounds}
\author{}
\date{8 October 2026}
\begin{document}
\maketitle

\begin{abstract}
This continuation records the mathematics developed after the morning
handoff of 8 October. The new material comprises an explicit rational
negative Einstein strictly almost-K\"ahler metric and its end geometry;
an exact harmonic-flux decomposition of traceless Ricci energy, including
its equality structure; a constrained local variational formulation and
actual descending deformations of products of negatively curved surfaces;
and global upper bounds for the relative Einstein defect, with an exact
unit-level calculation on a topological range. The constructions and
identities are proved in the text. Conditional local estimates state their
remaining geometric hypotheses. Earlier Burns insertion and gauge-theoretic
gap theorems enter only where needed as explicit inputs. The distinction
between an improved competitor, a local variation, and an unrestricted
global minimizer is retained throughout.
\end{abstract}

\paragraph{Relation to the preceding manuscript.}
The baseline is the pair
\path{01_Research_Synthesis.tex} and
\path{02_Mathematical_Source_Companion.tex}, whose mathematical cutoff was
8 October 2026, 09:58 CST. This file treats A01 Steps 56--58 and A02
Steps 65--72, including the completed Weyl-deficit argument and the
corrected proof versions.
It is independently compilable and contains no external input files.
Source identifiers are resolved in Appendix~\ref{app:sources}.
``Subsequent'' refers to this research interval, not to a claim that every
identity below is new to the literature.

\tableofcontents

\section{Objects, conventions, and the new connections}
\label{sec:conventions}

Unless stated otherwise, $M$ is a smooth closed connected oriented
four-manifold, and all infima are over all smooth Riemannian metrics.
Set
\begin{align}
 P(g)&=\int_M|\Ric^\circ|^2\dd,&
 S(g)&=\int_Ms^2\dd,&
 R(g)&=P(g)+\frac14S(g),\label{eq:energies}\\
 V(g)&=\Vol_g(M),& T_s(g)&=\int_Ms\dd,&
 L_s(g)&=S(g)-\frac{T_s(g)^2}{V(g)}.
\end{align}
The norm on two-forms is the form norm, so
$|e^{12}+e^{34}|^2=2$. The Weyl and full-curvature norms are curvature-operator
norms. Thus, with $T_\pm=2\chi(M)\pm3\tau(M)$,
\begin{equation}
 4\pi^2T_\pm=\frac{S}{24}+2W_\pm-\frac P2,
 \qquad W_\pm=\int_M|W^\pm|^2\dd,
 \qquad \int_M|\Rm_{\rm op}|^2\dd=8\pi^2\chi(M)+P.
 \label{eq:cw}
\end{equation}
Use different notation for the two variational problems:
\begin{equation}
 p(g)=\frac{P(g)}{8\pi^2},\qquad
 \pstar(M)=\inf_g p(g),\qquad
 P_*(M)=8\pi^2\pstar(M),
 \label{eq:pstar}
\end{equation}
and
\begin{equation}
 \cE(g)=\inf_{\lambda\ne0}
 \frac{\int_M|\Ric-\lambda g|^2\dd}{\lambda^2 V},
 \qquad \alpha(M)=\inf_g\cE(g).
 \label{eq:alpha}
\end{equation}
For $R>0$,
\begin{equation}
 \cJ(g)=\frac{4P}{R},\qquad
 \cE(g)=4-\frac{T_s^2}{VR}
        =\cJ(g)+\frac{L_s}{R}.
 \label{eq:EJ}
\end{equation}
In particular $\cE=\cJ$ for constant scalar curvature. If $T_s\ne0$,
the optimizing constant is $\lambda=R/T_s$. If $T_s=0$, the infimum
in~\eqref{eq:alpha} is $4$; for a Ricci-flat metric it is also $4$.
The exclusion of $\lambda=0$ is therefore mathematically consequential.

The developments fit into four connections:
\[
\begin{array}{c}
 \text{actual almost-K\"ahler equations}
 \longrightarrow \text{explicit complete noncompact geometry},\\[2pt]
 \text{harmonic Chern pairing}
 \longrightarrow \text{exact Ricci-energy decomposition and equality},\\[2pt]
 \text{equality geometry}
 \longrightarrow \text{local Hessians and genuine descending metrics},\\[2pt]
 \text{global competitors and topology}
 \longrightarrow \text{relative-gap upper bounds and exact unit levels}.
\end{array}
\]
These connections sharpen different parts of the same problem. None
identifies a local almost-K\"ahler solution with a closed Goldberg
counterexample, or a descending family with a global minimizer.

\section{An intrinsic equation and an explicit negative Einstein model}
\label{sec:ake}

\subsection{The actual almost-K\"ahler system}

Let $(g,J,\Omega)$ be almost-K\"ahler, with $|\Omega|^2=2$.
The Levi-Civita connection on $\Lambda^+$ has self-dual curvature precisely
when $g$ is Einstein. In an adapted frame $(\Omega,\phi,J\phi)$,
write its connection vector as $(b,Ja,a)$ and its curvature as
\begin{equation}
 F_1=db-a\wedge Ja=-\mathcal R(\Omega),\quad
 F_2=d(Ja)+a\wedge b=-\mathcal R(\phi),\quad
 F_3=da-Ja\wedge b=-\mathcal R(J\phi).
 \label{eq:frame-curvature}
\end{equation}
Here $\mathcal R$ denotes the curvature operator; these are the adapted-frame
Ricci identities of~\cite{AAD}. Put $\sigma=\Omega/\sqrt2$ and let
$\mathscr L=\sigma^\perp\subset\Lambda^+$ be the associated complex line.
The off-diagonal connection form $\zeta\in\Omega^1(\mathscr L)$ satisfies
\begin{equation}
 (d_b^*+d_b^-)\zeta=(2\eta,0),\qquad
 \zeta\in\Omega_J^{0,1}(\mathscr L),
 \label{eq:soldered}
\end{equation}
where the normalization of $\eta$ is fixed by
\[
 \nabla^*\nabla\Omega=2\eta+\tfrac12|\nabla\Omega|^2\Omega.
\]
The complex structure in the type condition is determined by the very
section being differentiated. This coupling distinguishes
\eqref{eq:soldered} from an arbitrary linear instanton deformation equation.

The canonical line connection is not, in general, a sourceless Maxwell
field. The first component of Bianchi is
\begin{equation}
 d^*F_1=\iota_{Ja}F_3-\iota_aF_2=:j.
 \label{eq:sourced}
\end{equation}
For constant $s$, the closed form
$\beta=-F_1-a\wedge Ja-(s/4)\Omega$ therefore satisfies
\[
 d^*\beta=-j-d^*(a\wedge Ja).
\]
In the Einstein almost-K\"ahler normalization used in \src{A01.56},
write $\beta=H+d\xi$, with $H$ harmonic, and put $q=|\nabla\Omega|^2$.
The self-dual/anti-self-dual components give
\begin{equation}
 \|d\xi\|_2^2=\frac1{16}\int q^2\dd-2\|H^-\|_2^2
 =\frac1{16}\int q^2\dd-8\pi^2\|c_1^{-,h}\|_2^2.
 \label{eq:hodge-remainder}
\end{equation}
Thus a Hodge decomposition does not remove the current in
\eqref{eq:sourced} or produce a uniform inverse estimate for $d+d^*$.

There is also a useful distinction between two differentiated tests.
Testing the transverse Bianchi identity against $\nabla_{JX}\Omega$
pairs its Weyl contribution with the skew Hessian, hence with curvature.
Testing against $\nabla_{h(X)}\Omega$, for harmonic $h\in\Omega^-$,
instead leaves
\begin{equation}
 \int\!\langle W_3^+,(\operatorname{Sym}_0\nabla^2\Omega)\circ h\rangle
 \quad\text{and}\quad
 \int\!\langle W_3^+,\nabla\Omega\circ\nabla h\rangle.
 \label{eq:unclosed-components}
\end{equation}
The second term sees the twistor component of $\nabla h$, even though
$dh=d^*h=0$. The first is not fixed by a commutator formula alone.
Consequently these tests identify concrete additional derivative information,
not an automatic sign for the harmonic twist pairing.

\subsection{Rational Bianchi VIII construction}
\label{sec:explicit}

Choose invariant one-forms on $\mathrm{SL}(2,\R)$ satisfying
\begin{equation}
 d\sigma_1=-2\sigma_2\wedge\sigma_3,\quad
 d\sigma_2=2\sigma_3\wedge\sigma_1,\quad
 d\sigma_3=2\sigma_1\wedge\sigma_2.
 \label{eq:VIII}
\end{equation}
For $1/3<x<1$, put $D=1+3x^2$ and define
\begin{align}
 G&=\frac{4D}{(1-x)^2(1+x)^2(3x-1)(3x+1)},\notag\\
 A_1&=\frac{4(3x-1)D}{3(1-x)(1+x)^2},&
 A_2&=\frac{4(3x+1)D}{3(1-x)^2(1+x)},&
 A_3&=\frac{64x^2}{(1-x)(1+x)D},\label{eq:VIII-metric}\\
 X_1&=\frac{(1-x)^3(3x+1)}{16x^3},&
 X_2&=\frac{(1+x)^3(3x-1)}{16x^3}.\label{eq:VIII-X}
\end{align}
Let
\[
 g=G\,dx^2+A_1\sigma_1^2+A_2\sigma_2^2+A_3\sigma_3^2,
 \quad e^0=\sqrt G\,dx,\quad e^i=\sqrt{A_i}\sigma_i,
\]
and $\omega_i=e^{0i}+e^{jk}$ for cyclic $(i,j,k)$.

\begin{proposition}[An actual strictly almost-K\"ahler Einstein pair]
\label{prop:VIII}
The form
\begin{equation}
 \Omega=\sqrt{X_1}\,\omega_1+\sqrt{X_2}\,\omega_2
 \label{eq:VIII-Omega}
\end{equation}
obeys
\begin{equation}
 d\Omega=0,\qquad |\Omega|^2=2,\qquad
 \Ric_g=-3g,\qquad W^-=0,
 \label{eq:VIII-properties}
\end{equation}
and
\begin{equation}
 q=|\nabla\Omega|^2
 =\frac{36(1-x^2)^3}{(1+3x^2)^3}>0.
 \label{eq:VIII-q}
\end{equation}
\end{proposition}

\begin{proof}
Here is a rational verification which does not require evaluating square
roots numerically. Write $dt=\sqrt G\,dx$, $v=dx/dt$, and
$h_i=(\log\sqrt{A_i})^{\boldsymbol\cdot}=vH_i$, where a dot denotes
$d/dt$ and a prime below denotes $d/dx$. Then
\begin{align*}
 v^2&=\frac{(1-x^2)^2(9x^2-1)}{4D},\qquad
 \ell=(\log v)'=-\frac{2x}{1-x^2}+\frac{9x}{9x^2-1}-\frac{3x}{D},\\
 H_1&=\frac12\left(\frac3{3x-1}+\frac{6x}{D}
               +\frac1{1-x}-\frac2{1+x}\right),\\
 H_2&=\frac12\left(\frac3{3x+1}+\frac{6x}{D}
               +\frac2{1-x}-\frac1{1+x}\right),\\
 H_3&=\frac1x+\frac{x}{1-x^2}-\frac{3x}{D}.
\end{align*}
Define the signed orbit coefficients
$w_i=n_i\sqrt{A_i/(A_jA_k)}=vV_i$, where
$(n_1,n_2,n_3)=(-1,1,1)$. Explicitly,
\begin{equation}
 V_1=-\frac{D}{4x(1+x)(3x+1)},\quad
 V_2=\frac{D}{4x(1-x)(3x-1)},\quad
 V_3=\frac{12x}{(9x^2-1)D}.
 \label{eq:VIII-V}
\end{equation}
The diagonal Ricci equations and their radial constraint reduce to
\begin{align}
 v^2\left[H_i'+\ell H_i+H_i\sum_jH_j
              -2\{V_i^2-(V_j-V_k)^2\}\right]&=3,\label{eq:VIII-ricci-cert}\\
 v^2\left[\sum_{i<j}H_iH_j+\sum_iV_i^2
                         -2\sum_{i<j}V_iV_j\right]&=3.\label{eq:VIII-ham-cert}
\end{align}
Substitution and a common denominator make each identity exact. Equivalently,
$\Ric_{00}=-\sum_i(\dot h_i+h_i^2)$ and
$\Ric_{ii}=-\dot h_i-h_i\sum_jh_j+2(w_i^2-(w_j-w_k)^2)$ are all $-3$.

The identity $X_1+X_2=1$ proves the length assertion.
Closure reduces, for $i=1,2$, to
\begin{equation}
 \mu_i:=2w_i-h_j-h_k=\frac v2(\log X_i)'.
 \label{eq:VIII-closure-cert}
\end{equation}
Let $C_i=w_i-w_j-w_k$ and $b_i=h_i+C_i$. Direct covariant
differentiation gives
\[
 q=2\{X_1\mu_1^2+X_2\mu_2^2+b_3^2+X_1b_2^2+X_2b_1^2\},
\]
which simplifies to~\eqref{eq:VIII-q}. Finally,
\begin{equation}
 (h_i-C_i)^{\boldsymbol\cdot}+h_i(h_i-C_i)=1
 \quad (i=1,2,3)
 \label{eq:VIII-Wminus-cert}
\end{equation}
is again a rational identity. Together with the Einstein equations it says
that $\mathcal R|_{\Lambda^-}=-\Id=(s/12)\Id$, proving $W^-=0$.
All denominators used above are nonzero on the specified interval.
\end{proof}

\subsection{End geometry and the role of compactness}
\label{sec:ends}

Let $r$ be geodesic distance from the $x=1/3$ end. Expanding the exact
radial equation gives
\begin{align}
 x-\tfrac13&=\tfrac29r^2-\tfrac2{81}r^4+O(r^6),\notag\\
 A_1&=r^2+\tfrac29r^4+O(r^6),&
 A_2&=6+7r^2+O(r^4),& A_3&=6+7r^2+O(r^4),\notag\\
 A_2-A_3&=\tfrac{10}{3}r^4+O(r^6),&
 X_2&=\tfrac83r^2+O(r^4).\label{eq:bolt-series}
\end{align}
The square roots and all coefficient quotients in these expressions are
analytic in $r^2$. Smoothness uses the compact circle in
$\mathrm{SL}(2,\R)$ with period $2\pi$. Its adjoint action on
$(\sigma_2,\sigma_3)$ has weight two. A period-$\pi$ quotient with the
same metric coefficients instead has cone angle $\pi$ at the collapsed
circle.

Locally write
$\sigma_1=d\psi+\beta$ and
$\sigma_2+i\sigma_3=e^{2i\psi}\zeta_0$, and set $z=re^{i\psi}$.
The combinations $dr\pm ir\sigma_1$ are expressed by
$dz+iz\beta$ and its conjugate with weight-one factors. In the metric,
the weight-four anisotropy $\sigma_2^2-\sigma_3^2$ is multiplied by
$r^4$ times a smooth even function. To check the mixed part of $\Omega$
explicitly, write $a_i=\sqrt{A_i}$, $a_1=r b(r^2)$ and $x_i=\sqrt{X_i}$.
It is
\[
 c_+\operatorname{Re}((dr+ir\sigma_1)\wedge(\sigma_2+i\sigma_3))
 +c_-\operatorname{Re}((dr-ir\sigma_1)\wedge(\sigma_2+i\sigma_3)),
 \qquad c_\pm=\tfrac12x_2(a_2\pm b a_3).
\]
Here $c_+=rF(r^2)$ and $c_-=r^3H(r^2)$, so their angular weights
become $z$ and $z^3$.
The radial--circle and horizontal area terms extend as well. These are
precisely the polar smoothness conditions; both $g$ and $\Omega$ extend
over the disk-bundle zero section.

Take a torsion-free cocompact surface-group lift
$\Gamma\subset\mathrm{SL}(2,\R)$ not containing $-I$.
The quotient principal orbits have the required $2\pi$ circle, so this
construction gives a smooth bolt over the compact hyperbolic surface.
At the other end, with $y=1-x$,
\begin{equation}
 G\sim\frac1{2y^2},\qquad
 A_1\sim\frac8{3y},\qquad
 A_2\sim\frac{32}{3y^2},\qquad
 A_3\sim\frac8y.
 \label{eq:infinite-end}
\end{equation}
The radial distance diverges, so the completed metric is complete and
noncompact. Its volume is infinite. More precisely,
\begin{equation}
 dV=\frac{64x(1+3x^2)}{3(1-x^2)^3}\,dx\wedge\sigma_1\wedge\sigma_2\wedge\sigma_3.
 \label{eq:VIII-volume}
\end{equation}
If $V_\Gamma=\int_{\Gamma\backslash\mathrm{SL}(2,\R)}\sigma_{123}$,
then direct one-variable integration gives
\begin{equation}
 \int q\,dV=64V_\Gamma,\qquad
 \int q^2\,dV=\frac{512}{3}V_\Gamma.
 \label{eq:VIII-integrals}
\end{equation}
For the above lift of a genus-$g$ surface,
$V_\Gamma=2\pi^2(g-1)$ in normalization~\eqref{eq:VIII}.
Thus $64$ and $512/3$ are radial integrals per unit invariant orbit volume.

The model has $q>0$ at every finite point and $q\to0$ at infinity.
It cannot descend, as this metric--form pair, to a compact quotient:
a positive continuous function on such a quotient would have a positive
minimum. It supplies an actual noncompact solution against which putative
pointwise negative-scalar integrability estimates must be tested.
Comparison with Armstrong's ansatz~\cite{ArmstrongAnsatz} is needed before
ascribing literature novelty to the family.

\begin{proposition}[A zero required in every closed candidate]
\label{prop:zero}
On a closed almost-K\"ahler four-manifold with $\Ric=\lambda g$,
$\lambda\ne0$, the Nijenhuis tensor vanishes somewhere. Equivalently,
$|\nabla\Omega|$ vanishes somewhere.
\end{proposition}
\begin{proof}
A nowhere-zero Nijenhuis tensor gives
$5\chi+6\tau=0$ by the characteristic-class calculation in
\cite{ArmstrongThesis}: it is a section of
$\Lambda^{0,2}T^*M\otimes T^{1,0}M$, whose second Chern number is
$c_2+2c_1^2=5\chi+6\tau$. In the almost-K\"ahler orientation this implies
$2\chi+3\tau=-\chi/2$. But Einstein Chern--Weil gives
$\chi>0$ when $s\ne0$, and~\eqref{eq:cw} gives
$2\chi+3\tau\ge0$. This is a contradiction.
\end{proof}
The topological input is classical. Its relevance here is to isolate what
the explicit solution lacks: a closed candidate must contain a torsion zero.
No propagation theorem proving that one such zero forces
$\nabla\Omega\equiv0$ has been obtained.

\section{Harmonic flux and the exact decomposition of Ricci energy}
\label{sec:flux}

\subsection{A metric quantity detected by every compatible field}

Write the curvature operator as
\begin{equation}
 \mathcal R=
 \begin{pmatrix}W^++s/12&B^*\\B&W^-+s/12\end{pmatrix},
 \qquad B:\Lambda^+\longrightarrow\Lambda^-,
 \qquad |\Ric^\circ|^2=4|B|^2.
 \label{eq:block}
\end{equation}
Let $a\in H^2(M;\R)$ and let $h_a^-$ be the anti-self-dual part of
its harmonic representative. Fix a nonzero harmonic anti-self-dual form $h$
and suppose
\[
 H=\|h\|_2^2,\qquad q_h=(h_a^-,h)>0.
\]
Parentheses are $L^2$ pairings: $(h_a^-,h)=-a\cdot[h]$, rather than
the intersection pairing itself.

For a $g$-orthogonal orientation-compatible almost-complex structure $J$
with $c_1(J)=a$, its canonical frame-curvature form obeys
\begin{equation}
 \gamma_J=\mathcal R(\Omega)+t_J,\qquad
 [\gamma_J]=2\pi a.
 \label{eq:canonical}
\end{equation}
Here $t_J$ is the quadratic intrinsic-torsion term. For example, in a
local adapted frame with
$\nabla\Omega=\alpha\otimes\phi+\beta\otimes J\phi$,
the signs can be fixed so that $t_J=\alpha\wedge\beta$.
Taking the anti-self-dual part of~\eqref{eq:canonical} and pairing gives
\begin{equation}
 (t_J^-,h)=2\pi q_h-\int_M\langle\Omega,B^*h\rangle\dd.
 \label{eq:chern-flux}
\end{equation}
Define quantities depending only on $g,a,h$ by
\begin{equation}
 \mathfrak F_h=\sqrt2\int_M|B^*h|\dd,\qquad
 \delta_h=1-\frac{\mathfrak F_h}{2\pi q_h},\qquad
 L_h=\frac{P/2-\mathfrak F_h^2/H}{4\pi^2}.
 \label{eq:flux-def}
\end{equation}

\begin{proposition}[Exact flux decomposition]
\label{prop:flux}
One has
\begin{equation}
 \boxed{p=\frac{q_h^2}{H}(1-\delta_h)^2+L_h,\qquad L_h\ge0,}
 \label{eq:flux-decomposition}
\end{equation}
and every $J$ in~\eqref{eq:canonical} satisfies
\begin{equation}
 (t_J^-,h)\ge2\pi q_h\delta_h.
 \label{eq:twist-lower}
\end{equation}
Off the zero set $Z(h)$, put
\[
 \beta_h=\frac{B^*h}{|h|^2},\quad
 B_\parallel(v)=\langle\beta_h,v\rangle h,\quad
 B_\perp=B-B_\parallel,\quad
 c=H^{-1}\int_M|\beta_h||h|^2\dd.
\]
Then
\begin{equation}
 \boxed{L_h=\frac1{2\pi^2}\left[
 \|B_\perp\|_2^2+
 \int_M(|\beta_h|-c)^2|h|^2\dd\right].}
 \label{eq:rank-amplitude}
\end{equation}
\end{proposition}
\begin{proof}
The norm bound $|\Omega|=\sqrt2$ proves~\eqref{eq:twist-lower}.
Cauchy--Schwarz yields
\[
 \mathfrak F_h^2
 \le2\left(\int|B||h|\dd\right)^2
 \le2\|B\|_2^2H=\frac P2H.
\]
Substitution proves~\eqref{eq:flux-decomposition}.
Pointwise orthogonal projection gives
$|B|^2=|B_\perp|^2+|\beta_h|^2|h|^2$.
The weighted variance identity
\[
 \int|\beta_h|^2|h|^2-
 H^{-1}\left(\int|\beta_h||h|^2\right)^2
 =\int(|\beta_h|-c)^2|h|^2
\]
then proves~\eqref{eq:rank-amplitude}. The zero set has measure zero,
so the displayed integrals are unchanged by its omission.
\end{proof}

For the unit form $\widehat h=h/\|h\|_2$ and
$q=(h_a^-,\widehat h)>0$, abbreviate $\delta=\delta_{\widehat h}$,
$L=L_{\widehat h}$. The formulas become
\begin{equation}
 p=q^2(1-\delta)^2+L,\quad
 \mathfrak F_{\widehat h}\le2\pi\sqrt p,\quad
 p<q^2\Longrightarrow\delta>0.
 \label{eq:unit-flux}
\end{equation}
Thus a saving below the squared harmonic projection requires a positive
metric flux deficit. The nonnegative remainder separates loss of rank
alignment from variation of the aligned amplitude.

The direction
$\Omega_{\rm Ric}=\sqrt2 B^*\widehat h/|B^*\widehat h|$
pointwise saturates the flux bound where defined. It need not extend to
a smooth global structure of Chern class $a$. Formula~\eqref{eq:twist-lower}
is valid for all class-$a$ structures without assuming this extension or
identifying an infimum over such structures with the pointwise bound.

\subsection{Equality and removable harmonic-form zeros}
\label{sec:flux-equality}

We regard two-forms as skew endomorphisms when writing their composition.
For a self-dual form $\phi$ and closed anti-self-dual form $h$,
\begin{equation}
 \diver(\phi\circ h)=h(d^*\phi),\qquad
 h^2=-\tfrac12|h|^2\Id.
 \label{eq:div-product}
\end{equation}
For example, in normal coordinates the derivative falling on $h$ in
$\nabla^a(\phi_{ac}h_{cb})$ vanishes by $dh=0$ and opposite duality;
the remaining term is $h_{bc}(d^*\phi)_c$.

\begin{theorem}[Constant-scalar equality structure]
\label{thm:flux-equality}
Assume $s$ is constant. If $L_h=0$, either $g$ is Einstein or it admits a
smooth compatible almost-K\"ahler structure $J_0$ with $J_0$-invariant
Ricci tensor and primitive Ricci form $\rho_0$ a constant multiple of $h$.
Conversely, any such almost-K\"ahler structure has $L_h=0$.
\end{theorem}
\begin{proof}
Equality in~\eqref{eq:rank-amplitude} implies
\[
 B(v)=c\langle\widehat\Omega_0,v\rangle h,
 \qquad |\widehat\Omega_0|=1,
\]
off $Z(h)$, with $c\ge0$ constant. If $c=0$, then $B=0$.
Otherwise let $\Omega_0=\sqrt2\widehat\Omega_0$.
The block--tensor correspondence is
\[
 B(v)=\langle\beta_h,v\rangle h
 \quad\Longrightarrow\quad
 \Ric^\circ=-2\beta_h\circ h=c'\Omega_0\circ h,
 \qquad c'=-\sqrt2c.
\]
Contracted Bianchi and~\eqref{eq:div-product} give
\[
 c'h(d^*\Omega_0)=\tfrac14ds=0.
\]
Since $h$ is invertible off its zero set, $d^*\Omega_0=0$ there; self-duality
also gives $d\Omega_0=0$.

The nodal-set theorem for the Hodge--Dirac operator~\cite{Baer} gives
$\dim_{\rm H}Z(h)\le2$, hence $\mathcal H^3(Z(h))=0$.
Covering this set by balls with arbitrarily small sum of cubed radii
produces cutoffs $\chi_j$ with $\|d\chi_j\|_{L^1}\to0$.
Because $|\Omega_0|=\sqrt2$, the cutoff boundary terms in its weak closed
and coclosed equations vanish. Elliptic regularity extends $\Omega_0$
smoothly across $Z(h)$. Its constant nonzero length persists.
The endomorphism product shows $J_0$-invariance of Ricci, and
$\rho_0=B\Omega_0=\sqrt2c\,h$ is harmonic.

Conversely, for an almost-K\"ahler structure the identity
$d\rho=*(Jd^*\Ric'')$ of~\cite{AAD} implies $d\rho_0=0$ when
$s$ is constant and $\Ric''=0$. Here $\Ric''$ is the $J$-anti-invariant
part of Ricci. The curvature block is
$B(v)=\tfrac12\langle\Omega,v\rangle\rho_0$.
If $\rho_0$ is a constant multiple of $h$, both terms
of~\eqref{eq:rank-amplitude} vanish.
\end{proof}

The theorem describes an equality set. It does not construct a deformation
joining a K\"ahler member to an Einstein member. Its strictly almost-K\"ahler
part is a constrained instance of the negative-scalar Blair--Ianus
integrability question, with prescribed harmonic alignment.

\subsection{The weaker rank-aligned system}

If only $B_\perp=0$, put $\phi_h=-2\beta_h$. Then
$\Ric^\circ=\phi_h\circ h$, and constant scalar curvature implies
$d\phi_h=d^*\phi_h=0$ off $Z(h)$.
If $\phi_h$ is locally bounded near $Z(h)$, the preceding removability
argument applies. The global two-form $F=\phi_h+h$ then satisfies
\begin{equation}
 dF=d^*F=0,\qquad
 \Ric^\circ=\tfrac12(F\circ F)_0.
 \label{eq:EM}
\end{equation}
Indeed, squares of self-dual or anti-self-dual skew endomorphisms are scalar,
and opposite-duality endomorphisms commute. Equivalently
$\widetilde F=(\phi_h-h)/\sqrt2$ obeys
$(\Ric+\widetilde F\circ\widetilde F)_0=0$.
This specifies the Einstein--Maxwell normalization.
Rank alignment alone controls $\int|\phi_h|^2|h|^2$, which does not
establish the boundedness needed at zeros. Full equality $L_h=0$ supplies
that boundedness through its constant amplitude.

For approximate alignment write
$\Ric^\circ=\phi_h\circ h+R_\perp$. If $s$ is constant, then for test
one-forms supported away from $Z(h)$,
\begin{equation}
 |(d^*\phi_h,\xi)|\le
 \|R_\perp\|_2\|\nabla(h^{-1}\xi)\|_2,
 \qquad \|R_\perp\|_2^2=8\pi^2L_{\rm rank}.
 \label{eq:weak-alignment}
\end{equation}
This follows by integrating
$h(d^*\phi_h)=-\diver R_\perp$.
It is a weighted weak estimate. The inverse $h^{-1}$ explains precisely
why small $L$ does not yet imply an unweighted global closeness theorem.

\subsection{Signed torsion and the scalar-curvature sign}
\label{sec:torsion}

On the almost-K\"ahler equality branch, take $\|\widehat h\|_2=1$ and
write $\rho_0=f\widehat h$, with $f$ constant and signed. Set
\[
 b=c_1(J),\qquad D_t=(t^-,\widehat h),\qquad
 \Phi_b=2\pi(h_b^-,\widehat h)=f+D_t.
\]
Then $P=2f^2$. If $\Phi_b\ne0$, define $d_b=D_t/\Phi_b$.
Unlike $\delta$ in~\eqref{eq:flux-def}, this definition uses signed flux:
\begin{equation}
 p=\frac{\Phi_b^2}{4\pi^2}(1-d_b)^2.
 \label{eq:signed-deficit}
\end{equation}
The coefficient equals $|b^-|^2$ only when
$\widehat h$ is parallel to $h_b^-$.

Let $u=|\nabla\Omega|^2/4$. In the present form norms,
$|t^-|^2=u^2/2$. The almost-K\"ahler Chern--Weil identity
\cite{AAD}, specialized to $\Ric''=0$, is
\[
 P=2\|\gamma^-\|_2^2+
 \int u(5u+s)\dd+2\|\rho_*''\|_2^2+4\|W_3^+\|_2^2.
\]
Here $\rho_*''$ is the anti-invariant star-Ricci component and $W_3^+$
is the tracefree Weyl block on $\Omega^\perp$. Substituting
$\gamma^-=\rho_0+t^-$ gives the exact balance
\begin{equation}
 \boxed{4fD_t+\int u(6u+s)\dd
       +2\|\rho_*''\|_2^2+4\|W_3^+\|_2^2=0.}
 \label{eq:torsion-balance}
\end{equation}
This calculation retains the sign of the harmonic pairing rather than
replacing it by the norm of the twist.

If $s\ge0$ and $\Phi_b>0$, Cauchy--Schwarz gives
$D_t^2\le\tfrac12\int u^2$, whence
\[
 4fD_t\le-6\int u^2\le-12D_t^2.
\]
Using $f=\Phi_b-D_t$ yields
\begin{equation}
 -\tfrac12\le d_b\le0,\qquad
 p\ge\frac{\Phi_b^2}{4\pi^2}.
 \label{eq:positive-s-twist}
\end{equation}
Here $d_b=0$ forces $u=0$ and hence K\"ahlerness. For constant $s<0$,
the same computation instead gives
\begin{equation}
 \boxed{-s\int u\dd\ge4\Phi_b^2d_b(1+2d_b).}
 \label{eq:negative-s-twist}
\end{equation}
In particular, if $\widehat h\parallel h_b^-$ and
$p\le(1-\eta)^2|b^-|^2$ for $0<\eta<1$, then
$\eta\le d_b\le2-\eta$ and
\[
 -s\int u\dd\ge4\Phi_b^2\eta(1+2\eta).
\]
Thus a prescribed saving requires an actual amount of nonintegrability.
At negative scalar curvature, $d_b=0$ alone no longer implies integrability.

\subsection{What the low-cost cell measures}

The already constructed perpendicular cell can be used to evaluate the
decomposition, rather than only its total energy. For the retained cell
with parameters $x_0=0.03$ and $X=128$, floating quadrature gives
\[
 p\simeq0.995663,\quad H\simeq0.99897,\quad
 \delta_{\rm metric}\simeq0.00541,\quad
 H^{-1}(1-\delta_{\rm metric})^2\simeq0.99023,\quad
 L\simeq0.00543.
\]
Here the relative pairing is normalized to $q_h=1$. The exact bookkeeping
identity is
\[
 1-p=2\delta-\delta^2-L-(1-\delta)^2(H^{-1}-1).
\]
These are numerical diagnostics on a previously certified cell, not
interval bounds for the flux variables. They show why the energy decrease
cannot be inferred from the twist alone: its contribution is offset by
misalignment and the harmonic normalization. The geometric identities
above, rather than these particular numbers, are the reusable result.

\section{Local variational geometry at a constant-scalar K\"ahler metric}
\label{sec:variation}

\subsection{Stationarity of the corrected functional}

Let $(g_0,J_0,\omega)$ be constant-scalar K\"ahler, with
$a=c_1(J_0)$ and $\rho_0=2\pi h_a^-\ne0$.
For nearby metrics take
$\widehat h_g=h_{a,g}^-/\|h_{a,g}^-\|_2$ and
$q_g=\|h_{a,g}^-\|_2$. At $g_0$ the flux decomposition gives
$\delta=L=0$. The natural local functional is
\begin{equation}
 \mathscr D_a(g)=p(g)-|a_g^-|^2
 =C_+(g)-|a_g^+|^2,
 \quad
 C_+(g)=\frac1{2\pi^2}\int|W^+|^2\dd+
                 \frac1{96\pi^2}\int s^2\dd.
 \label{eq:corrected-functional}
\end{equation}
The second identity follows from~\eqref{eq:cw} and
$|a^+|^2-|a^-|^2=a^2=2\chi+3\tau$ for this almost-complex class.

\begin{proposition}[Stationarity under unrestricted metric variations]
\label{prop:stationarity}
At a constant-scalar K\"ahler metric,
\begin{equation}
 D\delta=DL=0,\qquad D\mathscr D_a=0.
 \label{eq:stationarity}
\end{equation}
The statement about $\mathscr D_a$ also holds in the K\"ahler--Einstein
case, where the normalized deficit itself is undefined.
\end{proposition}
\begin{proof}
Along any two-sided smooth metric variation $g_t$, orthogonalize $J_0$
to obtain a compatible $J_t$ in the same homotopy class. Its intrinsic
torsion is $O(t)$, so the quadratic form $t_{J_t}$ in
\eqref{eq:canonical} is $O(t^2)$. Inequality~\eqref{eq:twist-lower}
therefore gives $\delta(g_t)\le C t^2$.
Differentiability at zero and both signs of $t$ give $D\delta=0$.
The norm integral is differentiable here even if $h_a^-$ has zeros:
its zero set has measure zero and the derivative is dominated by the
smooth variation of $B^*\widehat h$.
Nonnegativity and $L(g_0)=0$ give $DL=0$.
Finally
\[
 \mathscr D_a=L-2q_g^2\delta+q_g^2\delta^2
\]
proves stationarity. At a K\"ahler--Einstein metric,
$P=|a^-|^2=0$, and both have zero first derivative separately.
\end{proof}
This is stationarity of $p-|a^-|^2$, not stationarity of $P$ at every
constant-scalar K\"ahler metric.

\subsection{The period obstruction in a local coercivity argument}

The following formulation separates the established linear analysis from
the cohomological input which it still needs. Work on the linear space
\[
 \mathcal S=\{k:\beta_{g_0}k=0,\quad s'_0(k)
                    \text{ is spatially constant}\},
\]
where $\beta$ is Bianchi gauge. Let $d_2(k)$ and $\ell_2(k)$ be the
coefficients of $t^2$ in $\delta(g_0+tk)$ and $L(g_0+tk)$.
They are coefficients, not second derivatives. Then
\begin{equation}
 (\mathscr D_a)_2(k)=\ell_2(k)-2q_0^2d_2(k),\qquad
 c_K=\inf_{d_2(k)>0}\frac{\ell_2(k)}{q_0^2d_2(k)}.
 \label{eq:cK}
\end{equation}

An $\ell_2$-null direction follows the moving harmonic line, not a fixed
background form. Its alignment equation includes
\[
 B'_0=\dot\Phi\otimes h_0+\Phi_0\otimes\dot h
\]
with the metric-induced bundle identifications understood. Linearized
Bianchi and the amplitude condition then give an aligned pair to first
order:
\begin{equation}
 \Ric_t^\circ=\Phi_t\circ h_t+O(t^2),\quad
 d\Omega_t=O(t^2),\quad (\Ric'')'_0=0.
 \label{eq:first-order-pair}
\end{equation}
To compare this pair with a K\"ahler variation one must also match the
cohomology class of the varying closed form. Define
\begin{equation}
 \mathcal P_\omega:H^{0,1}(T^{1,0}M)\longrightarrow H^{0,2}(M),
 \quad [\mu]\longmapsto[\mu\mathbin{\lrcorner}\omega],\qquad
 \Xi(k)=[\dot\Omega_k]^{0,2}\pmod{\operatorname{im}\mathcal P_\omega}.
 \label{eq:Xi}
\end{equation}
Vanishing of $\Xi(k)$ supplies a first-order K\"ahler comparison with
the same class variation. After subtracting that comparison, a Moser
correction removes the exact form variation. In the form norms used here,
the fixed-symplectic Blair--Ianus identities are
\begin{equation}
 E_\omega=\int|\nabla\omega|^2\dd,\qquad
 \int s\dd=4\pi c_1\cdot[\omega]-\tfrac12E_\omega,\qquad
 DE_\omega(k)=2\int\langle\Ric'',k\rangle\dd.
 \label{eq:blair}
\end{equation}
At a K\"ahler metric,
$\operatorname{Hess}E_\omega(k,k)=2\|(\nabla\omega)'_k\|_2^2$.
Equation~\eqref{eq:first-order-pair} therefore makes a null direction with
$\Xi(k)=0$ first-order K\"ahler; it lies in the radical of $d_2$.

\begin{proposition}[Conditional local coercivity]
\label{prop:conditional-local}
Suppose $h_{a,g_0}^-$ is zero-free. Put
$K=\ker\ell_2\cap\mathcal S$. The augmented alignment, scalar, and
Bianchi system is elliptic and $K$ is finite-dimensional. If
$\Xi|_K=0$, then $c_K>0$, and for every $0<c<c_K$,
\begin{equation}
 L(g_0+k)\ge c q_0^2\delta(g_0+k)
                   -C\|k\|_{C^{2,\gamma}}^3,
 \qquad k\in\mathcal S\text{ sufficiently small},\quad0<\gamma<1.
 \label{eq:local-coercivity}
\end{equation}
\end{proposition}
\begin{proof}
On the Bianchi-symbol kernel the Ricci symbol is
$\tfrac12|\xi|^2k$. Vanishing of the augmented symbol gives
\[
 k^\circ=\sigma\circ\widehat h,quad \sigma\perp\omega,\quad
 \tr k=0,\quad (\sigma\circ\widehat h)\xi=0.
\]
Opposite duality implies
$(\sigma\circ\widehat h)^2=\tfrac14|\sigma|^2|\widehat h|^2\Id$.
Since $\widehat h$ is zero-free, the symbol is injective. The resulting
elliptic estimate on the closed manifold is
\[
 \ell_2(k)\ge c_1\|k-\pi_Kk\|_{H^2}^2.
\]
The first-order K\"ahler comparisons just described put $K$ in the
radical of $d_2$. The positive part of the torsion quadratic form is
first order, giving
\[
 d_2(k)\le C\|k-\pi_Kk\|_{H^1}^2\le C'\ell_2(k).
\]
Taylor expansion of $L$ and $\delta$ proves~\eqref{eq:local-coercivity}.
\end{proof}
The required conclusion $\Xi|_K=0$ has not been proved in general.
Nor is $\mathcal S$ itself a nonlinear constant-scalar slice; use of
\eqref{eq:local-coercivity} on that locus requires its parametrization
and the quadratic correction. Along a ray with $d_2>0$, a cubic remainder
is generally of order $\delta^{3/2}$, not $\delta^2$.

One further formula records the moving-frame term in this analysis.
For $v_t=B_t^*\widehat h_t$, $e_t=\Omega_t^{\rm ref}/\sqrt2$ and
$w_1=P_{e_0^\perp}v'_0-|v_0|e'_0$,
\begin{equation}
 2\pi q_0d_2=\tau_{{\rm ref},2}
             -\frac1{\sqrt2}\int\frac{|w_1|^2}{|v_0|}\dd.
 \label{eq:delta-Hessian}
\end{equation}
Here $\tau_{{\rm ref},2}$ is the quadratic coefficient of the reference
torsion pairing. On $K^\perp$, in the $\ell_2$ Hilbert norm, this is an
operator of the form $T-N$, with $T$ compact self-adjoint and $N\ge0$
bounded. Its positive spectrum is discrete; compactness of the entire
operator is neither needed nor supplied.

\subsection{A genuine descending family on products of curves}
\label{sec:product}

The local constant cannot be uniform over degenerating backgrounds, as
an explicit calculation shows. Let
\[
 M=\Sigma_1(K_1)\times\Sigma_2(K_2),\qquad K_1,K_2<0,quad K_1\ne K_2,
\]
with constant-curvature factor metrics. Let $Q\ne0$ be the real tensor
of a holomorphic quadratic differential on $\Sigma_1$, and take
$\Delta_2f=\mu f$, $\mu>0$, with nonnegative Laplacian. Set
$A=g_1^{-1}Q$ and consider the actual metrics
\begin{equation}
 g_t=g_1(e^{tfA}\,\cdot,\cdot)\oplus g_2.
 \label{eq:product-family}
\end{equation}
Since $A$ is tracefree, both area forms remain fixed. Consequently
$\omega_1-\omega_2$ remains harmonic anti-self-dual and
$q_t^2=|c_1^-|_{g_t}^2$ is constant. The tangent $k=fQ$ is transverse
traceless and $s'_0=0$.

\begin{proposition}[Product Hessian]
\label{prop:product}
The Taylor coefficient of order two is
\begin{equation}
 \boxed{8\pi^2(\mathscr D_a)_2
      =\frac\mu4(\mu+K_1-K_2)\|fQ\|_2^2.}
 \label{eq:product-Hessian}
\end{equation}
If $K_1<K_2$ and $0<\mu<K_2-K_1$, this gives strict local descent.
The descent persists after making scalar curvature exactly constant.
\end{proposition}
\begin{proof}
Let $G_t$ denote the first metric block and $K_t$ its Gaussian curvature
at a fixed point of the second factor. The block connection, through the
terms required at order two, gives
\[
 \Ric_V=K_tG_t+\frac{t\mu f}{2}G_tA,
 \qquad
 \Ric_H=K_2g_2-\frac{t^2}{4}|Q|^2df\otimes df,
\]
and mixed Ricci has no term of order one. With
$n_t=t^2|Q|^2|df|^2/4$, this implies
\[
 |\Ric^\circ|^2=(K_t-K_2)^2
       +\frac{t^2\mu^2f^2}{4}|Q|^2+(K_t-K_2)n_t+O(t^3).
\]
The transverse-traceless property gives $K'_0=0$. For each second-factor
point, Gauss--Bonnet and the fixed area form give
$\int_{\Sigma_1}(K_t-K_1)dA_1=0$. Hence the integrated order-two
contribution from $(K_t-K_2)^2$ is zero. Finally
$\int|df|^2=\mu\int f^2$ yields~\eqref{eq:product-Hessian}.

Since $s_0<0$, the linearized conformal scalar operator
$6\Delta-2s_0$ is invertible. An $O(t^2)$ conformal correction therefore
makes scalar curvature constant. Stationarity of $\mathscr D_a$ leaves
its second coefficient unchanged. Hodge decomposition on two-forms is
conformally invariant, so the conclusion also gives the stated decrease
in $P$ along the corrected family.
\end{proof}

The separate coefficients of the flux decomposition are
\begin{equation}
 \ell_2(k)=\frac{\mu^2}{32\pi^2}\|k\|_2^2,
 \qquad
 q_0^2d_2(k)=\frac{(K_2-K_1)\mu}{64\pi^2}\|k\|_2^2.
 \label{eq:product-split}
\end{equation}
For example, the form-norm calculation gives
$|(\nabla\omega)'|^2=|Q|^2|df|^2/2$ and
$\int u_2=\mu\|k\|_2^2/8$, consistent with these constants.
Thus
\begin{equation}
 \boxed{c_K\le\frac{2\lambda_1(\Sigma_2)}{K_2-K_1}.}
 \label{eq:spectral-degeneration}
\end{equation}
Pinching a separating geodesic produces $\lambda_1(\Sigma_2)\to0$:
test with a function nearly constant on the two thick components and
varying across their increasingly long collar. The gradient energy tends
to zero while its mean-zero $L^2$ norm stays positive. Therefore there is
no positive uniform lower bound for $c_K$ even at fixed product topology
and fixed unequal curvature scales.

The modes in~\eqref{eq:product-split} have $\ell_2>0$, so they do not
decide $\Xi$ on the nullspace. For products of curves of genus at least
two, $\mathcal P_\omega=0$: there are no cross-factor holomorphic vector
fields, and the factor deformations preserve the type of their area
classes. It remains to determine whether the constrained nullspace has
a nonzero $\Xi$ component. Also, descent below $|c_1^-|^2$ on this
product is not a construction below the blow-up cost $k$: the latter
problem has the additional projection contribution in
$|c_1^-|^2=k+\Pi$.

\section{From actual competitors to global normalized upper bounds}
\label{sec:upper}

\subsection{The optimized Yamabe conversion}

Let $Y([g])$ be the Yamabe constant, and $\widehat g$ a minimizing
Yamabe representative. In a negative Yamabe class,
\begin{equation}
 S(\widehat g)=Y([g])^2,
 \qquad P(g)-P(\widehat g)=\tfrac1{12}(S(g)-Y([g])^2)\ge0.
 \label{eq:Yamabe-P}
\end{equation}
The equality uses conformal invariance of Weyl energy and
\eqref{eq:cw}. For a direct positive-remainder formula, write
$g=v^2\widehat g$, $s_{\widehat g}=s_0<0$. With all integrals and
derivatives on the right taken in $\widehat g$,
\begin{equation}
 S(g)-Y([g])^2
 =36\int\left(\frac{\Delta v}{v}\right)^2dV_{\widehat g}
       -12s_0\int\frac{|\nabla v|^2}{v^2}dV_{\widehat g}.
 \label{eq:Yamabe-exact}
\end{equation}
This follows by squaring
$s_g=v^{-3}(6\Delta v+s_0v)$ and integrating by parts.

\begin{proposition}[Conversion of a $P$ competitor]
\label{prop:conversion}
For every $g$ with $Y([g])<0$,
\begin{equation}
 \boxed{\alpha(M)\le
 \frac{16P(g)}{4P(g)+Y([g])^2}.}
 \label{eq:P-to-alpha}
\end{equation}
If every conformal class has negative Yamabe constant and
$Y([g])^2\ge32\pi^2A$ with the same $A>0$, then
\begin{equation}
 \boxed{\alpha(M)\le\frac{4\pstar(M)}{A+\pstar(M)}.}
 \label{eq:global-P-upper}
\end{equation}
For either orientation with $T_\pm>0$, one also has the unconditional
lower comparison
\begin{equation}
 \boxed{\alpha(M)\ge
 \frac{\pstar(M)}{\pstar(M)+3T_\pm/4}.}
 \label{eq:global-P-lower}
\end{equation}
\end{proposition}
\begin{proof}
At $\widehat g$, minimizing in $\lambda$ gives
\[
 \cE(\widehat g)=\frac{16P(\widehat g)}{4P(\widehat g)+Y([g])^2}.
\]
Monotonicity in $P$, followed by~\eqref{eq:Yamabe-P}, proves
\eqref{eq:P-to-alpha}; take a $P$-minimizing sequence for
\eqref{eq:global-P-upper}. For the lower bound,~\eqref{eq:cw} implies
\[
 R=4P+24\pi^2T_\pm-12W_\pm\le4P+24\pi^2T_\pm.
\]
Combine this with $\cE\ge4P/R$ and monotonicity in $P$.
If a Ricci-flat metric is present, it has $\cE=4$, which also satisfies
the proposed lower bound.
\end{proof}

The numerator $16P/Y^2$ is only the trial value at $\lambda=s/4$,
not the optimized expression. The denominator in
\eqref{eq:P-to-alpha} is essential even for a constant-scalar metric.
On the stated supply manifolds with $T_+>0$, the two comparisons imply
\begin{equation}
 \alpha(M)=0\quad\Longleftrightarrow\quad P_*(M)=0.
 \label{eq:conditional-zero-equivalence}
\end{equation}
They do not assert equality of the two functionals or a sharp small-gap
coefficient.

For
$M=X\#k\overline{\CP}^{\,2}\#\ell(S^1\times S^3)$,
where $X$ is minimal of general type and $A=c_1^2(X)>0$,
the polarized scalar supply in LeBrun's Proposition 3.2~\cite{LeBrun}
provides the input to~\eqref{eq:global-P-upper}. The relevant statement
includes handles and the appropriate $b^+=1$ polarization; one must not
replace it by the false assertion that every general-type surface has
$b^+>1$.

\subsection{The new consequence of the earlier Burns construction}

The morning manuscript already contains the following construction input:
for every minimal general-type surface $X$ and every $k\ge1$,
\begin{equation}
 \text{there is a smooth metric on }X\#k\overline{\CP}^{\,2}
 \text{ with }P<8\pi^2k.
 \label{eq:inherited-Burns}
\end{equation}
On a negative K\"ahler--Einstein host this comes from the perpendicular
Burns insertion, whose normalized cost at $s=-4$ is
\[
 1-\tfrac14\varepsilon^2\log(1/\varepsilon)+O(\varepsilon^2).
\]
If the canonical model has ADE points, the saving is fixed first and the
ALE resolution scales are then made sufficiently small. This order is
part of the inherited construction, not a new result here.

Applying Proposition~\ref{prop:conversion} gives its strengthened
normalized consequence:
\begin{equation}
 \boxed{\alpha(X\#k\overline{\CP}^{\,2})
               <\frac{4k}{A+k}.}
 \label{eq:strict-general-type-upper}
\end{equation}
The new point is the passage to the optimized relative functional without
separate scalar-variance or volume estimates.

\subsection{A nonasymptotic law on all further finite covers}
\label{sec:cover-law}

The retained interval certificate for the complex-hyperbolic blow-up cell
has upper endpoint
\begin{equation}
 p_c=\frac{338806471362324583565914367157197246841}{2^{128}}
       =0.995662732771114200\ldots<1.
 \label{eq:pc}
\end{equation}
For every $\epsilon>0$ it supplies a smooth cell of topology
$D^4\#\overline{\CP}^{\,2}$, exactly complex hyperbolic near its
outer boundary, with $P/(8\pi^2)\le p_c+\epsilon$.
Fix its support radius $R_c$. A handle cell with the same outer collar
has costs approaching $P=32\pi^2$.

Let $Y=\Gamma\backslash\mathbb B^2$ be a compact ball quotient normalized
by $\Ric=-g$. Choose a compact fundamental set $F$, $x_0\in\mathbb B^2$,
and $D_F=\sup_{x\in F}d(x_0,x)$. Residual finiteness yields a normal
finite-index subgroup excluding all nonidentity elements satisfying
\[
 d(x_0,\gamma x_0)\le4R_c+2D_F.
\]
If an element of this subgroup had displacement at most $4R_c$ at any
point, conjugate that point into $F$; normality and the triangle inequality
would contradict this exclusion. The resulting cover $X'$ thus has
$\operatorname{inj}(X')>2R_c$ everywhere. Every finite cover $X\to X'$
inherits the same bound.

A maximal $2R_c$-separated set in $X$ has disjoint $R_c$-balls and its
$2R_c$-balls cover $X$. Since $\Vol(X)=2\pi^2A$, $A=c_1^2(X)$, the
number $N(X)$ of such centers satisfies
\begin{equation}
 N(X)\ge\frac{2\pi^2A}{v_{\rm CH}(2R_c)}.
 \label{eq:packing}
\end{equation}

\begin{theorem}[Finite-cover upper law]
\label{thm:cover-law}
For every finite cover $X\to X'$ and every integer $n\ge1$ with
$15n\le N(X)$, set
\[
 M=X\#8n\overline{\CP}^{\,2}\#7n(S^1\times S^3),\quad
 Q=36n,\quad h=Q/A.
\]
Then
\begin{equation}
 \boxed{\frac{\alpha(M)}h\le
  \frac{C_c}{1+C_ch/4},\qquad
 C_c=\frac{28}{9}+\frac89p_c
     =3.9961446513521015108\ldots<4.}
 \label{eq:cover-upper}
\end{equation}
The topology satisfies $T_+=A-Q$ and $T_-=(A+Q)/3$.
\end{theorem}
\begin{proof}
Insert $8n$ blow-up and $7n$ handle cells in the chosen disjoint balls.
The complement is Einstein, so locality of the $P$ integral and arbitrarily
small smoothing errors give
\[
 \pstar(M)\le8np_c+28n=\frac{C_cQ}{4}.
\]
Substitute into~\eqref{eq:global-P-upper}. The displayed topological
formulas follow from $T_+(X)=A$, $T_-(X)=A/3$, a change $(-1,+5)$
per blow-up, and $(-4,-4)$ per handle.
\end{proof}
There is one preliminary cover, followed by universal quantifiers over
all further finite covers and all admissible $n$. Neither nesting nor an
$h\to0$ limit is needed. The estimate concerns an infimum, so it does
not claim that one smoothed metric attains the interval endpoint.

The complete handle insertion has a useful scale calculation. Take the
scalar-flat conformally flat asymptotically Euclidean Green-function
metric on $(S^1\times S^3)\setminus\{p\}$. Its single-end
Gauss--Bonnet formula gives $P=32\pi^2$. Insert a copy scaled by $\delta$
and match at radius $t=\sqrt\delta$. The background normal-coordinate
error is $O(t^2)$ and the scaled AE error is $O(\delta^2t^{-2})$.
After two derivatives and integration over the annulus,
\begin{equation}
 P_{\rm annulus}=O(t^4+\delta^4t^{-4})=O(\delta^2).
 \label{eq:handle-error}
\end{equation}
The omitted tail contributes $O((\delta/t)^4)$ and the removed background
ball contributes $O(t^4)$. Thus on any fixed smooth background
\[
 P(\widetilde g)=P(g)+32\pi^2+O(\delta^2).
\]
All constants here may depend on that background; the estimate is not a
uniform gluing theorem for arbitrary degenerating backgrounds.

\subsection{Fixed negative K\"ahler--Einstein hosts}

On a fixed negative K\"ahler--Einstein surface, choose finitely many
disjoint normal-coordinate balls and a fixed sufficiently small Burns
parameter giving saving $\eta_X>0$. Each lifted ball in any finite
Riemannian cover has the same jets and the same error constants.
For $k$ blow-ups and $\ell$ handles within this packing capacity, put
$Q=k+4\ell$, $h=Q/A$, and $\theta=k/Q$. The same argument gives
\begin{equation}
 \frac{\alpha(M)}h\le
 \frac{4(1-\theta\eta_X)}{1+h(1-\theta\eta_X)}.
 \label{eq:KE-cover-upper}
\end{equation}
This statement applies to covers that exist; it does not posit an
unbounded tower on a simply connected surface.

It is also different from balanced small-$h$ asymptotics.
If the fixed host is not a ball quotient, then
$\Delta_X=3\chi(X)-c_1^2(X)>0$; under unramified covers,
$\Delta/A=\Delta_X/A_X$ stays fixed. In the earlier balance parameter
\[
 \sigma=\frac{4\Delta}{Q}+18\theta
        =\frac{4\Delta_X}{hA_X}+18\theta,
\]
one has $\sigma\to\infty$ as $h\to0$.
The upper construction still applies, but the fixed balanced small-$h$
window does not follow from these towers. Varying hosts with
$\Delta/A\to0$ require the separate ramified-cover construction.

\section{The unit level and a global conformal variational formula}
\label{sec:weyl-law}

\subsection{A smooth long-cylinder competitor on every manifold}

\begin{theorem}[Universal upper bound]
\label{thm:universal}
For every closed connected smooth four-manifold,
\begin{equation}
                    0\le\alpha(M)\le1.
 \label{eq:universal}
\end{equation}
\end{theorem}
\begin{proof}
Choose a fixed smooth rotationally symmetric cap $C\cong D^4$ with
metric $dr^2+f(r)^2g_{S^3(1)}$, where $f=\sin r$ near the center and
$f=1$ near the boundary. A smooth cutoff constant to all orders at its
ends gives such a cap. Join two copies by $[0,L]\times S^3(1)$.
The result is a smooth metric on $S^4$. Replace a fixed ball in one cap
by a fixed punctured copy of $M$ through a fixed smooth neck. This is
$M\#S^4\cong M$. The complement $\Omega$ of the length-$L$ cylinder
is literally independent of $L$.

On the cylinder,
\[
 \Ric=(0,2,2,2),\quad s=6,\quad
 |\Ric|^2=12,\quad |\Ric^\circ|^2=3.
\]
Writing $a_3=\Vol(S^3)=2\pi^2$, the integrals therefore have the form
\[
 V=a_3L+v_0,\quad T_s=6a_3L+t_0,\quad
 R=12a_3L+r_0,\quad P=3a_3L+p_0.
\]
Put $b=\int_\Omega(s-6)dV$ and $J_0=\int_\Omega(s-6)^2dV$.
Then exactly
\[
 \overline s=6+b/V,\qquad L_s=J_0-b^2/V=O(1).
\]
Thus $\cE(g_L)=(4P+L_s)/R\to1$. There is no change of topology and
no growing contribution from a transition region.
\end{proof}

The special long-circle construction on
$N\#(S^1\times Y^3)$, with $Y$ spherical or hyperbolic, uses the same
fixed-complement calculation. The spherical dumbbell shows that a handle
summand is unnecessary. This is a statement about the relative functional:
$P(g_L)$ itself tends to infinity.

\subsection{Exact values at and beyond the topological wall}

Combining~\eqref{eq:EJ} and~\eqref{eq:cw}, for $R>0$,
\begin{equation}
 \cE(g)-1=
 \frac{12W_\pm+L_s-24\pi^2T_\pm}{R}.
 \label{eq:unit-wall}
\end{equation}
Let $T_{\min}=2\chi-3|\tau|$.

\begin{corollary}[Value and attainment]
\label{cor:unit-wall}
If $T_{\min}\le0$, then $\alpha(M)=1$.
If $T_{\min}<0$, the infimum is not attained.
If $T_{\min}=0$, an attaining metric is exactly a half-conformally-flat
metric of nonzero constant scalar curvature, with the vanishing Weyl half
selected by the orientation for which $T_\pm=0$.
\end{corollary}
\begin{proof}
All terms in the numerator of~\eqref{eq:unit-wall} are nonnegative
in the stated orientation. Theorem~\ref{thm:universal} gives equality
at the infimum. If $T_{\min}<0$, every $R>0$ metric satisfies
\[
 \cE(g)\ge1+\frac{24\pi^2|T_{\min}|}{R}>1.
\]
A Ricci-flat metric has $\cE=4$, so it cannot attain $1$ either.
If $T_{\min}=0$, equality requires $W_\pm=0$ and $L_s=0$.
The constant scalar curvature cannot be zero:~\eqref{eq:cw} would then
give $P=0$, and the original definition would give $\cE=4$.
The converse follows directly from~\eqref{eq:unit-wall}.
\end{proof}

For $T_{\min}<0$, every minimizing sequence necessarily has
$R\to\infty$, $W_\pm=o(R)$, and $L_s=o(R)$.
For $T_{\min}=0$, attainment is equivalently the existence of a
half-conformally-flat conformal class with nonzero Yamabe constant.
If $\tau=0$, the signature identity makes this full conformal flatness.
In the positive-scalar equality case, the harmonic-form Weitzenb\"ock
formula gives $b_+=0$, so $b_-=4-4b_1$; no conclusion of simple
connectivity follows from this Betti-number calculation.

Some resulting values and examples are
\begin{equation}
 \alpha(K3)=\alpha(T^4)=\alpha(S^2\times T^2)=1,
 \qquad \alpha(\#_n\overline{\CP}^{\,2})=1\quad(n\ge4).
 \label{eq:unit-examples}
\end{equation}
The first three values do not by themselves decide attainment.
Products $S^1\times(S^3/\Gamma)$ and orientation-preserving isometric
mapping tori of closed hyperbolic three-manifolds attain $1$ with their
locally product metrics. Their Euler characteristic is zero; an Einstein
metric would consequently be flat, contrary to their fundamental groups.
LeBrun's positive-Yamabe half-flat conformal classes on
$\#_4\overline{\CP}^{\,2}$~\cite{LNN} give another attained example.
These examples are Bach-flat. They do not supply the non-Bach-flat
global $P$-minimizer sought in the third public problem.

The parameter $h=(k+4\ell)/A$ alone does not characterize the unit level.
For a ball quotient with $k=0$ and $A/12<\ell<A/4$,
\[
 h<1,\qquad T_-=A/3-4\ell<0,
\]
so $\alpha=1$ nonetheless. Thus a converse to
``$h\ge1$ implies $\alpha=1$'' is false for mixed surgery families.

\subsection{The orientation-free Weyl deficit}
\label{sec:weyl-deficit}

Let
\begin{equation}
 \mathcal W([g])=\int_M|W|^2\dd=W_++W_-,\qquad
 d([g])=8\pi^2\chi(M)-\mathcal W([g]).
 \label{eq:d-def}
\end{equation}
Both are conformal invariants in dimension four. Gauss--Bonnet gives
\begin{equation}
 P=\frac S{12}-2d,\qquad R=\frac S3-2d,
 \qquad
 \boxed{\cE(g)=1+\frac{L_s(g)-6d([g])}{R(g)}}\quad(R>0).
 \label{eq:orientation-free}
\end{equation}
This removes the orientation from~\eqref{eq:unit-wall} and exhibits
the actual two-variable conformal problem.

\begin{lemma}[Conformal minimization below the unit level]
\label{lem:conformal-min}
Every representative of a conformal class satisfies $S(g)\ge Y([g])^2$.
If $d([g])>0$, then $Y([g])\ne0$, and a Yamabe representative
minimizes $\cE$ in that class. Its value is
\begin{equation}
 \cE(\widehat g)=1-\frac{6d([g])}{Y([g])^2/3-2d([g])}.
 \label{eq:class-value}
\end{equation}
\end{lemma}
\begin{proof}
If $Y>0$, the test function $1$ in the Yamabe quotient gives
$Y\le T_s/\sqrt V\le\sqrt S$.
If $Y\le0$, H\"older's inequality gives, for every nonzero test function,
\[
 \frac{\int(6|\nabla u|^2+su^2)\dd}{(\int u^4\dd)^{1/2}}
 \ge-\|s_-\|_2.
\]
Taking the infimum yields $|Y|\le\|s_-\|_2\le\sqrt S$.
At a Yamabe representative $P=Y^2/12-2d\ge0$.
Consequently $d>0$ implies $Y^2\ge24d>0$ and
$R(\widehat g)=Y^2/3-2d\ge6d>0$.
For every representative $R(g)\ge R(\widehat g)$ and
\[
 \cE(g)\ge1-\frac{6d}{R(g)}
          \ge1-\frac{6d}{R(\widehat g)}=\cE(\widehat g).
\]
Equality requires $L_s=0$ and $S=Y^2$.
\end{proof}

\begin{theorem}[Global conformal variational law]
\label{thm:global-law}
On every closed oriented four-manifold,
\begin{equation}
 \boxed{\alpha(M)=1-6\sup\left(
 \{0\}\cup\left\{
 \frac{d([g])}{Y([g])^2/3-2d([g])}:d([g])>0
 \right\}\right).}
 \label{eq:global-law}
\end{equation}
Each displayed positive ratio belongs to $(0,1/6]$.
Writing $\mathcal W_*(M)=\inf_g\int|W|^2\dd$, one has
\begin{equation}
 \boxed{\alpha(M)<1
       \quad\Longleftrightarrow\quad
       \mathcal W_*(M)<8\pi^2\chi(M).}
 \label{eq:weyl-threshold}
\end{equation}
\end{theorem}
\begin{proof}
Classes with $d\le0$ have $\cE\ge1$ by~\eqref{eq:orientation-free};
any Ricci-flat representative has value $4$.
Classes with $d>0$ are treated by Lemma~\ref{lem:conformal-min}.
The universal competitor in Theorem~\ref{thm:universal} supplies the
value $1$ as an infimum even if the set of positive-deficit classes is empty.
This proves~\eqref{eq:global-law} including that case.
A positive supremum is equivalent to the existence of one
positive-deficit class, which is~\eqref{eq:weyl-threshold}.
\end{proof}

This formula has no monopole or negative-Yamabe hypothesis. It is a
variational characterization, not a computation of its supremum on
every manifold. In particular the infimum of Weyl energy alone decides
whether $\alpha<1$, but it does not determine the amount below $1$;
the Yamabe value of the same conformal class enters the denominator.
The joint Yamabe--Weyl viewpoint has classical antecedents~\cite{ABKS}.

On the negative-Yamabe supply manifolds of
Proposition~\ref{prop:conversion}, put
$p_Y=P(\widehat g)/(8\pi^2)$ and
$\upsilon=Y([g])^2/(32\pi^2)-A$. Then
\begin{equation}
 \alpha(M)=\min\left\{1,
       \inf_{[g]}\frac{4p_Y}{p_Y+A+\upsilon}\right\}.
 \label{eq:supply-quotient}
\end{equation}
The exact conformal reduction below $1$ was already present in the
earlier manuscript. The new formula~\eqref{eq:global-law} removes that
supply restriction and includes the unit-level alternative.

\subsection{Connected sums and an exact Weyl threshold}

The standard small-neck construction for Weyl energy gives
\begin{equation}
 \mathcal W_*(M\#N)\le\mathcal W_*(M)+\mathcal W_*(N).
 \label{eq:W-subadd}
\end{equation}
One can see the estimate directly: make each fixed metric Euclidean in
a sufficiently small normal ball. Since $g_{ij}-\delta_{ij}=O(r^2)$,
this costs $o(1)$ Weyl energy. Join the flat collars by a conformally
flat rotationally symmetric neck, which has zero Weyl energy.
Take nearly minimizing metrics first and then send the ball radii to zero.

For $d_*(M)=8\pi^2\chi(M)-\mathcal W_*(M)$,
Euler-characteristic additivity now gives
\begin{equation}
 \boxed{d_*(M\#N)\ge d_*(M)+d_*(N)-16\pi^2.}
 \label{eq:deficit-sum}
\end{equation}
The neck loses $16\pi^2$ of Euler contribution. For $S^4$,
$d_*=16\pi^2$, so a sphere summand is neutral as it must be.
For $\overline{\CP}^{\,2}$, reversed Fubini--Study and the signature
bound give $\mathcal W_*=12\pi^2$, $d_*=12\pi^2$.
Induction and the same signature lower bound yield
\begin{equation}
 \mathcal W_*(\#_n\overline{\CP}^{\,2})=12n\pi^2,\qquad
 d_*(\#_n\overline{\CP}^{\,2})=4\pi^2(4-n).
 \label{eq:ncp-W}
\end{equation}
Hence, for $n\ge1$,
\begin{equation}
 \alpha(\#_n\overline{\CP}^{\,2})<1\quad(n\le3),\qquad
 \alpha(\#_n\overline{\CP}^{\,2})=1\quad(n\ge4).
 \label{eq:ncp-threshold}
\end{equation}
This threshold combines classical Weyl-energy information with the
relative-functional characterization; it is not a new construction of
the classical half-flat metrics.

For a smooth negative K\"ahler--Einstein surface $X$ with $A=c_1^2(X)$,
the Einstein metric has $d=4\pi^2A/3$. Consequently
\begin{equation}
 d_*(X\#k\overline{\CP}^{\,2})\ge4\pi^2(A/3-k).
 \label{eq:KE-deficit}
\end{equation}
Canonical models with ADE points require the separate orbifold/ALE
transfer, rather than this smooth-KE argument alone.
For any such blow-up, the identity $T_+=A-k$ and the signature formula
give the exact reformulation
\begin{equation}
 \boxed{\alpha(X\#k\overline{\CP}^{\,2})<1
 \quad\Longleftrightarrow\quad
 \inf_g\int|W_g^+|^2dV_g<2\pi^2(A-k).}
 \label{eq:blowup-W}
\end{equation}
Thus the intermediate range $A/3\le k<A$ contains a definite
all-conformal-class Weyl-energy question. The strict construction
\eqref{eq:strict-general-type-upper} already gives $\alpha<1$ for
$k\le A/3$ when that integer endpoint is present; it does not compute
the remaining values in this intermediate range.

\subsection{What a zero-infimum statement must control}

Formula~\eqref{eq:EJ} gives the exact criterion
\[
 \cE(g_j)\to0\quad\Longleftrightarrow\quad
 \frac{4P(g_j)+L_s(g_j)}{R(g_j)}\to0
\]
when $R(g_j)>0$. A Ricci-flat bulk with $P\to0$ can fail this criterion.
For example $P_*(K3)=P_*(T^4)=0$, whereas~\eqref{eq:unit-examples}
gives $\alpha=1$. Collapsing the flat factor of
$S^2\times T^2$ similarly gives $P_*=0$, not $\alpha=0$.

For a sequence with $\cE(g_j)\to0$, rescale by its nonzero optimizing
Einstein constant so that $\lambda_{\rm opt}=\pm1$. Then
\begin{equation}
 R=|T_s|=(4-\cE)V,
 \qquad
 48\pi^2\chi=6\mathcal W+L_s+(4-\cE)(1-\cE)V.
 \label{eq:normalized-small}
\end{equation}
The second equality follows by substituting the first in Gauss--Bonnet.
In particular $V\le12\pi^2\chi(1+o(1))$.
A bubble-free, globally bounded-curvature total collapse of this same
normalized sequence would force $\chi=0$, contradicting
\eqref{eq:normalized-small} for subunit metrics. This observation does
not transfer a compactness theorem for a specially selected
$P$-minimizing sequence to arbitrary $\cE$-minimizing sequences.

Likewise, an ALE gluing estimate for $P$ becomes a relative-gap rate only
after controlling $L_s$ and the nonzero bulk limit of $R$. In the
previous Eguchi--Hanson calculation the proposed $a^4/6$ coefficient
requires $L_s=o(a^4)$; the $P$-rate alone does not supply it.

\section{The mathematical questions now exposed}
\label{sec:questions}

The increment supplies an actual negative Einstein nonintegrable model,
an equality theory for the harmonic-flux remainder, descending closed
metrics, stronger global upper-bound quantifiers, and a general
Weyl--Yamabe expression for $\alpha$. Their relation to the four original
questions can be stated without introducing another research programme.

\begin{center}\small
\begin{tabularx}{\textwidth}{@{}p{.20\textwidth}XX@{}}
\toprule
Question & What the increment supplies & What remains to be proved\\
\midrule
Einstein existence on nonminimal general-type surfaces &
Exact flux/equality structure and a global conformal characterization &
Exclusion or construction of Einstein metrics in the remaining small-blow-up range\\
Optimal total curvature &
Product descent, stronger normalized recovery bounds, and exact relative-gap values in the unit region &
The unrestricted $P_*$, or matching asymptotics, on the target blow-up and balanced families\\
Non-Bach-flat global $P$-minimizer &
A more precise local variational problem, including the period obstruction &
A smooth attaining metric with the required non-Bach-flat property, or an exclusion theorem\\
Closed negative Goldberg &
A complete noncompact rational solution and the necessary closed torsion zero &
A compact integrability argument or a closed strictly almost-K\"ahler Einstein example\\
\bottomrule
\end{tabularx}
\end{center}

Four concrete connections would change these conclusions.

\begin{question}[The compact zero and differentiated compatibility]
Can the actual equations near a zero of the Nijenhuis tensor, together
with a global identity, force integrability? A successful estimate must
handle the two components in~\eqref{eq:unclosed-components} and remain
consistent with Proposition~\ref{prop:VIII}. A purely pointwise
negative-scalar sign argument cannot exclude that actual solution.
\end{question}

\begin{question}[Coercivity after alignment]
Can the weighted estimate~\eqref{eq:weak-alignment} be strengthened across
the harmonic zero set, and can $\Xi|_K$ be decided on the relevant
blow-ups? Equation~\eqref{eq:spectral-degeneration} shows that constants
uniform over arbitrary degenerating products cannot be obtained by
ignoring the small spectrum. The remaining information must come from
the target geometry or an additional global relation.
\end{question}

\begin{question}[A sharp joint conformal bound]
On $X\#k\overline{\CP}^{\,2}$, determine the supremum in
\eqref{eq:global-law}, rather than optimizing $P$ and $Y$ independently.
In particular, in the unresolved intermediate range, can one decide
\eqref{eq:blowup-W}? In a balanced small-$h$ family, can an all-metric
lower bound match a genuine recovery construction with coefficient below
$4$? The current cell coefficient is an upper bound, not a prescribed
answer.
\end{question}

\begin{question}[From a local critical structure to a global minimum]
Which compactness or direct global comparison turns the corrected
Hessian information into a statement about all competing metrics?
For nonzero-scale orbifold limits, the scalar-variance term must be
estimated at the same order as the traceless-Ricci energy. For the
non-Bach-flat problem, an attained relative gap at $1$ is not a substitute
for a global minimizer of $P$.
\end{question}

\appendix
\section{Mathematical source map}
\label{app:sources}

Project paths below are relative to the research archive;
\path{Sessions/} paths are relative to the Sessions' user workspace.
The manuscript does not require these files to compile.
\begingroup\small\raggedright
\begin{description}[style=nextline]
\item[Baseline]
\path{handoff/20261008_Einstein_Gaps_Four_Problems/01_Research_Synthesis.tex}
and \path{02_Mathematical_Source_Companion.tex} in the same directory.
The negative-background Burns saving and its ADE extension are inherited
from this pair. The older exact conformal comparison is also in
\path{skills/geometric-analysis/research/Curvature_Realizability.tex},
Proposition \texttt{vg-conformal-infimum}.
\item[A01.56]
\path{Sessions/S1003A01/workspace/notes/step56_soldered_instanton_and_SU2_AKE.md}.
Adapted-connection system, sourced current, Hodge remainder and the
two differentiated Bianchi tests. The symmetry constraint-tower calculations
lead into the following rational construction.
\item[A01.57--58]
\path{Sessions/S1003A01/workspace/notes/step57_bianchiVIII_strict_AKE_negative.md};
\path{Sessions/S1003A01/workspace/notes/step58_verified_bianchiVIII_AKE_and_closed_zero_locus.md}.
The complete rational and polar-extension proofs are in
\path{Sessions/S1003A01/trace/ttcs/results/step58_ttcs01_verifier1_integration.md}
and \path{step58_ttcs01_sub4_verifier.md} in that directory.
The exact calculation is reproduced by
\path{Sessions/S1003A01/workspace/scripts/step58_exact_bianchiVIII_fast.py}.
Equations~\eqref{eq:VIII-ricci-cert}--\eqref{eq:VIII-Wminus-cert}
retain the rational identities in the present standalone file.
\item[A02.65--66]
\path{Sessions/S1003A02/workspace/notes/step65_cell_exceptional_flux_ledger.md};
\path{Sessions/S1003A02/workspace/notes/step66_rank_one_exceptional_ricci_rigidity.md},
with its Section 7 corrections and
\path{Sessions/S1003A02/trace/ttcs/results/step66_ttcs01_verifier1_integration.md}.
These provide the flux decomposition, removable equality structure,
Einstein--Maxwell normalization and signed torsion balance.
\item[A02.67--68]
\path{Sessions/S1003A02/workspace/notes/step67_cK_local_coercivity_at_cscK.md};
\path{Sessions/S1003A02/workspace/notes/step68_product_cscK_saddle_verification.md}.
The final constants and period condition are in
\path{Sessions/S1003A02/trace/ttcs/results/step68_ttcs01_verifier1_integration.md}
and \path{step68_ttcs01_verifier2_integration.md} in that directory.
\item[A02.69]
\path{Sessions/S1003A02/workspace/notes/step69_P_form_upper_law_balanced_towers.md},
Section 6 and
\path{Sessions/S1003A02/trace/ttcs/results/step69_ttcs01_verifier1_integration.md}.
The inherited cell interval endpoint is in
\path{Sessions/S1003A01/workspace/data/step33/rigorous_x00.03_X128_d20_P_exact.json}.
\item[A02.71--72]
\path{Sessions/S1003A02/workspace/notes/step71_functional_split_and_unit_wall.md};
\path{Sessions/S1003A02/workspace/notes/step72_orientation_free_weyl_deficit_law.md},
including its final conformal-minimization argument in
\path{Sessions/S1003A02/trace/ttcs/results/step72_ttcs02_verifier1_integration.md}.
The smooth long-cylinder proof, both signs of the Yamabe comparison,
and the complete global formula are written out in
Section~\ref{sec:weyl-law} above.
\end{description}
\endgroup

\begin{thebibliography}{10}
\bibitem{AAD}
V. Apostolov, J. Armstrong and T. Dr\u aghici,
\emph{Local models and integrability of certain almost K\"ahler 4-manifolds},
Math. Ann. \textbf{323} (2002), 633--666.
\href{https://arxiv.org/abs/math/0109098}{arXiv:math/0109098}.
\bibitem{ArmstrongThesis}
J. Armstrong, \emph{Almost-K\"ahler Geometry}, doctoral thesis.
\href{https://johnarmstrongmaths.com/publications/resources/thesis.pdf}{Author's full text}.
The characteristic-class obstruction is used in Proposition~\ref{prop:zero}.
\bibitem{ArmstrongAnsatz}
J. Armstrong, \emph{An ansatz for almost-K\"ahler, Einstein 4-manifolds},
J. Reine Angew. Math. \textbf{542} (2002), 53--84.
\href{https://doi.org/10.1515/crll.2002.009}{doi:10.1515/crll.2002.009}.
\bibitem{Baer}
C. B\"ar, \emph{Zero sets of solutions to semilinear elliptic systems
of first order}, Invent. Math. \textbf{138} (1999), 183--202.
\href{https://arxiv.org/abs/math/9805065}{arXiv:math/9805065}.
\bibitem{LeBrun}
C. LeBrun, \emph{Ricci curvature, minimal volumes, and Seiberg--Witten theory}.
\href{https://arxiv.org/abs/math/0003068}{arXiv:math/0003068}.
Proposition 3.2 supplies the polarized scalar estimate used in
Section~\ref{sec:upper}.
\bibitem{Optimal}
C. LeBrun, \emph{Curvature Functionals, Optimal Metrics, and the
Differential Topology of 4-Manifolds}.
\href{https://arxiv.org/abs/math/0404251}{arXiv:math/0404251}.
\bibitem{LNN}
C. LeBrun, S. Nayatani and T. Nitta,
\emph{Self-dual manifolds with positive Ricci curvature}.
\href{https://arxiv.org/abs/dg-ga/9411001}{arXiv:dg-ga/9411001}.
\bibitem{ABKS}
K. Akutagawa, B. Botvinnik, O. Kobayashi and H. Seshadri,
\emph{The Weyl functional near the Yamabe invariant},
J. Geom. Anal. \textbf{13} (2003), 1--20.
\href{https://arxiv.org/abs/math/0201153}{arXiv:math/0201153}.
\end{thebibliography}
\end{document}
